\documentclass[11pt,letterpaper]{article}
\usepackage{jahnel}

\runningtitle{Subcritical isotropic Poisson line Cox percolation}
\hypersetup{pdftitle={On the subcritical phase of isotropic Poisson line Cox percolation},pdfauthor={Christian Hirsch, Benedikt Jahnel}}

\title{\textbf{On the subcritical phase of isotropic\\ Poisson line Cox percolation}\\[0.6em] \large Working note}

\author{%
  Christian Hirsch\thanks{Institut f\"ur Mathematische Stochastik, Georg-August-Universit\"at G\"ottingen, Goldschmidtstra\ss e~7, 37077 G\"ottingen, Germany. \texttt{christian.hirsch@uni-goettingen.de}. \orcid{0000-0003-4136-3740}.}
  \and
  Benedikt Jahnel\thanks{Institut f\"ur Mathematische Stochastik, Technische Universit\"at Braunschweig, Universit\"atsplatz~2, 38106 Braunschweig, Germany, and Weierstrass Institute for Applied Analysis and Stochastics, Mohrenstra\ss e~39, 10117 Berlin, Germany. \texttt{benedikt.jahnel@tu-braunschweig.de}. \orcid{0000-0002-4212-0065}.}%
}

\date{\today}

\begin{document}

\maketitle
\thispagestyle{plain}

\begin{abstract}
We collect what can presently be said about the subcritical phase of the planar Boolean model with disks centered at Poisson points on isotropic Poisson lines. We do not prove that the critical intensity is positive. Instead, we reduce this question to a statement about the line process alone. Given the lines, percolation is a zero--one event, and the resulting quenched threshold is almost surely constant and equal to the annealed one. A union bound over lattice paths then shows that the critical intensity is positive as soon as the exponential growth rate of path sums of line counts is finite with positive probability. We further show that connection probabilities never decay exponentially: for every intensity they are at least $\exp(-(4+o(1))(\log n)^2/(\lambda r))$ over distance~$n$, owing to bundles of nearly parallel lines. Annealed moments of the relevant path sums grow superexponentially, so the direct expectation bound does not close the gap. We discuss unresolved difficulties with decoupling, sprinkling and a vacant-circuit construction, and state the remaining greedy-path estimate as a conjecture.

\medskip
\noindent\textit{2020 Mathematics Subject Classification.} Primary 60K35; Secondary 60D05, 60G55.

\noindent\textit{Key words and phrases.} Continuum percolation, Poisson line Cox process, Boolean model, subcritical phase, greedy lattice animals.
\end{abstract}

% =================================================================
\section{Introduction}
\label{sec:introduction}

Let $X_\lambda$ be a Poisson line Cox process: the points of a homogeneous Poisson process of intensity~$\lambda$ on each line of an isotropic Poisson line process of intensity~$\mu$. The companion note~\cite{HJ26} shows that the Boolean model $\Xi_\lambda$ with radius~$r$ percolates once $\lambda\ge Cr^{-1}[1+\logp(1/(\mu r))]$. The complementary question, whether $\Xi_\lambda$ fails to percolate for small~$\lambda$, is raised by Hil\'ario, S\'a, Sanchis, and Teixeira~\cite[Section~9]{HSST26}. For the Manhattan environment of horizontal and vertical lines, Jahnel, Jhawar, and Vu~\cite{JJV23} establish phase transitions by discretization and comparison with percolation on randomly stretched lattices~\cite{Hof05}; this comparison has no isotropic counterpart. General subcritical criteria for Cox percolation~\cite{HJC19} require stabilization, which Poisson line length measures lack~\cite[Section~3]{JT22}.

The difficulty is visible in one example. A bundle of $j$ nearly parallel lines crossing a strip of length~$n$ and width of order~$r$ makes every point of the strip close to $j$ lines, so the strip carries about $j\lambda$ points per unit length. A bundle with $j$ of order $\log n/\lambda$ fills the strip with an occupied path, and such a bundle has probability $\exp(-O((\log n)^2/\lambda))$. Hence long occupied paths are never exponentially unlikely (\cref{prop:noexp}), and the argument for the supercritical phase in~\cite{HJ26}, which deletes directions, has the wrong monotonicity here~\cite[Remark~6.1]{HJ26}.

\paragraph{Results.} We prove four statements. First, percolation is monotone in all parameters, and the quenched percolation probability given the lines is $0$ or~$1$ with an almost surely constant threshold that equals the annealed critical intensity (\cref{prop:monotone,prop:quenched}). Second, connection probabilities decay at most like $\exp(-(4+o(1))(\log n)^2/(\lambda r))$ for every $\lambda>0$ (\cref{prop:noexp}). Third, and this is the main point, the critical intensity is positive as soon as a growth rate~$g^*$ defined from the line process alone is finite with positive probability (\cref{thm:criterion}). Fourth, the annealed moments of the corresponding path sums grow superexponentially (\cref{prop:annealed}), which rules out the direct expectation bound for these path sums.

\paragraph{Status.} We do not know a proof that $g^*<\infty$. \Cref{sec:greedy} proves a Borel--Cantelli bound on lines shared by distant boxes, formulates the missing estimate as a greedy-path conjecture, and explains the heuristic reasons to expect it. \Cref{sec:obstacles} records difficulties with several natural approaches. The companion file \texttt{corrected-geometry.pdf} revises the later greedy-path calculation and proves a bound for individual intervals along pairs of lines; the full greedy estimate remains open. All heuristic statements are labeled as such.

% =================================================================
\section{Setting}
\label{sec:setting}

We use the notation of~\cite{HJ26}. Put $n_\theta\coloneqq(\cos\theta,\sin\theta)$ and $\ell(\theta,s)\coloneqq\{x\in\R^2\colon x\cdot n_\theta=s\}$ for $\theta\in[0,\pi)$ and $s\in\R$. Let $\cL$ be a Poisson process of lines with intensity measure $\nu_\mu(\d\theta\,\d s)\coloneqq(\mu/\pi)\dd\theta\dd s$, let $\Lambda(A)\coloneqq\sum_{\ell\in\cL}\cH^1(\ell\cap A)$, and, conditionally on~$\cL$, let $X_\lambda$ be a Poisson process with intensity measure~$\lambda\Lambda$ \cite[Theorem~5.6]{LP17}. For $r>0$ put $\Xi_\lambda\coloneqq\bigcup_{x\in X_\lambda}B(x,r)$, where $B(x,t)$ is the open disk of radius~$t$ centered at~$x$. Percolation means the existence of an unbounded connected component of~$\Xi_\lambda$. Write $\theta(\lambda,\mu,r)$ for its probability and
\begin{equation*}
  \lambda_c(\mu,r)\coloneqq\inf\{\lambda>0\colon\theta(\lambda,\mu,r)>0\}\,.
\end{equation*}
For a convex body~$K$, the lines meeting~$K$ have $\nu_\mu$-measure $(\mu/\pi)$ times the perimeter of~$K$, by integrating the width of~$K$ over all directions. Two open disks of radius~$r$ intersect if and only if their centers are at distance less than~$2r$.

We repeatedly use lines that cross a strip lengthwise. Such lines meet every vertical slice of the strip, so a family of them puts the same amount of line length into every box of the strip.

\begin{lemma}[Lines crossing a strip]
\label{lem:strip}
Let $0<a\le d$ and $R\coloneqq[0,d]\times[0,a]$. The lines meeting both vertical sides of~$R$ have $\nu_\mu$-measure at least $\mu a^2/(5\pi d)$ and at most $2\mu a/\pi$. Each such line meets every slab $R\cap([t,t+w]\times\R)$, $0\le t\le d-w$, in a segment of length at least~$w$.
\end{lemma}

\begin{proof}
For the lower bound, consider the lines through $(0,y_0)$ with $y_0\in[a/4,3a/4]$ and slope $b\in[-a/(4d),a/(4d)]$. At $x=d$ their height $y_0+bd$ lies in $[0,a]$, so they meet both sides. Their normal angles fill an interval of length $2\arctan(a/(4d))\ge0.47\,a/d$ around~$\pi/2$, because $a/(4d)\le1/4$ and $\arctan t\ge0.94\,t$ for $0\le t\le1/4$. At a fixed normal angle~$\theta$, the displacement of the line through $(0,y_0)$ is $s=y_0\sin\theta$, which ranges over an interval of length $(a/2)\sin\theta\ge0.485\,a$, since $\sin\theta\ge(1+1/16)^{-1/2}$. Hence the measure is at least $(\mu/\pi)\cdot0.47\cdot0.485\,a^2/d\ge\mu a^2/(5\pi d)$. The upper bound is the measure of all lines meeting the left side, which is $(\mu/\pi)\cdot2a$. A line meeting both vertical sides is the graph of an affine function with values in $[0,a]$ at $x=0$ and $x=d$, hence on $[0,d]$; over $[t,t+w]$ it stays in~$R$ and has length at least~$w$.
\end{proof}

% =================================================================
\section{Monotonicity and the quenched threshold}
\label{sec:quenched}

The model is monotone in each parameter, and it has an exact scaling. Together these reduce the question to one dimensionless parameter and show that the critical intensity can only decrease when lines are added.

\begin{proposition}[Monotonicity and scaling]
\label{prop:monotone}
The percolation probability $\theta(\lambda,\mu,r)$ is nondecreasing in each of $\lambda,\mu,r$. Consequently, $\lambda_c(\mu,r)$ is nonincreasing in~$\mu$ and in~$r$. Moreover, $\lambda_c(\mu,r)=r^{-1}\lambda_c(\mu r,1)$.
\end{proposition}

\begin{proof}
For $\mu\le\mu'$, the superposition of independent Poisson line processes with intensities~$\mu$ and $\mu'-\mu$ is a Poisson line process with intensity~$\mu'$; marking all lines with independent Poisson processes of intensity~$\lambda$ couples the two Cox processes by inclusion. For $\lambda\le\lambda'$, independent thinning of the points with retention probability $\lambda/\lambda'$ gives the coupling. Increasing~$r$ enlarges every disk. Percolation is increasing under all three couplings. Under $x\mapsto x/r$, lines are mapped to lines, the intensity measure becomes $(\mu r/\pi)\dd\theta\dd s$, the linear point intensity becomes~$\lambda r$, and the radius becomes~$1$; this gives the scaling identity.
\end{proof}

The next lemma is the mixing property of the line process used in~\cite[Lemma~3.4]{HJ26}, stated for events of the lines alone. For a Borel set~$W$ write $[W]$ for the set of lines meeting~$W$.

\begin{lemma}[Mixing of the line process]
\label{lem:mixing}
Let $W_1,W_2$ be bounded Borel sets, and let $F,H$ be events determined by the lines in $[W_1]$ and $[W_2]$, respectively. Then $\lvert\P(F\cap H)-\P(F)\P(H)\rvert\le3\nu_\mu([W_1]\cap[W_2])$. If $W_1,W_2\subseteq B(0,\rho)$, then $\nu_\mu([W_1]\cap[W_2+x])\le4\mu\rho^2/\lvert x\rvert$ for $\lvert x\rvert\ge4\rho$. In particular, $\cL$ is mixing and hence ergodic under translations.
\end{lemma}

\begin{proof}
The restrictions of~$\cL$ to $[W_1]\setminus[W_2]$, to $[W_2]\setminus[W_1]$, and to $[W_1]\cap[W_2]$ are independent Poisson processes. On the event that the last restriction is empty, $F$ and~$H$ are determined by the first two restrictions respectively, so they are conditionally independent there. Since this event has probability at least $1-\nu_\mu([W_1]\cap[W_2])$, comparing the joint probability and then the two marginal factors with the corresponding probabilities on this event gives the bound. A line meeting both $B(0,\rho)$ and $B(x,\rho)$ has normal angle with $\lvert x\cdot n_\theta\rvert\le2\rho$, a set of angles of measure $2\arcsin(2\rho/\lvert x\rvert)\le2\pi\rho/\lvert x\rvert$ because $\arcsin t\le\pi t/2$ on $[0,1]$, and for each such angle its displacement lies in an interval of length~$2\rho$. Hence the measure is at most $(\mu/\pi)(2\pi\rho/\lvert x\rvert)(2\rho)=4\mu\rho^2/\lvert x\rvert$. Mixing for events determined by bounded windows extends to all events by approximation, and mixing implies ergodicity.
\end{proof}

Conditionally on~$\cL$, the model is a Poisson Boolean model with the locally finite intensity measure~$\lambda\Lambda$. The following proposition shows that conditioning on the lines loses nothing: it suffices to find, for almost every line configuration, some positive intensity without percolation, and this intensity may depend on the configuration.

\begin{proposition}[Quenched zero--one law and threshold]
\label{prop:quenched}
For almost every realization of~$\cL$ and every $\lambda>0$, the conditional percolation probability $\theta_\cL(\lambda)\coloneqq\P(\Xi_\lambda\text{ percolates}\,\vert\,\cL)$ lies in $\{0,1\}$ and is nondecreasing in~$\lambda$. The quenched threshold $\lambda_c(\cL)\coloneqq\sup\{\lambda>0\colon\theta_\cL(\lambda)=0\}$ is almost surely equal to the constant $\lambda_c(\mu,r)$. In particular, $\lambda_c(\mu,r)>0$ as soon as $\P(\lambda_c(\cL)>0)>0$.
\end{proposition}

\begin{proof}
Fix~$\cL$ such that only finitely many lines meet every bounded set; this holds almost surely. Then $X_\lambda$ is a Poisson process with locally finite intensity measure, and percolation does not change if the points in a bounded set~$B$ are modified. Indeed, adding points keeps an unbounded component inside an unbounded component. Removing the finitely many points in~$B$ from an unbounded component removes finitely many vertices from its connected and locally finite intersection graph; every remaining component contains a neighbor of a removed vertex, so there are finitely many of them, one of them is infinite, and its union of disks is unbounded because the points are locally finite. Since the restrictions of~$X_\lambda$ to disjoint annuli are independent, Kolmogorov's zero--one law gives $\theta_\cL(\lambda)\in\{0,1\}$. Independent thinning gives monotonicity in~$\lambda$.

The function $\cL\mapsto\lambda_c(\cL)$ is measurable, since it suffices to take the supremum over rational~$\lambda$, and it is invariant under translations of~$\cL$, because a translated line configuration produces the translated Poisson Boolean model and percolation is translation invariant. By \cref{lem:mixing} it is almost surely equal to a constant~$c^*$. For $\lambda<c^*$ we have $\theta_\cL(\lambda)=0$ almost surely, so $\theta(\lambda,\mu,r)=\E\theta_\cL(\lambda)=0$. For $\lambda>c^*$ there is a rational $\lambda'<\lambda$ with $\theta_\cL(\lambda')>0$, hence $\theta_\cL(\lambda)=1$ almost surely, and $\theta(\lambda,\mu,r)=1$. Thus $c^*=\lambda_c(\mu,r)$.
\end{proof}

In particular, $\theta(\lambda,\mu,r)\in\{0,1\}$ for every $\lambda\neq\lambda_c(\mu,r)$, in agreement with the almost-sure percolation in~\cite[Theorem~2.1]{HJ26}.

% =================================================================
\section{No exponential decay}
\label{sec:noexp}

In a subcritical Poisson Boolean model with bounded radii, connection probabilities decay exponentially. Here they cannot: a bundle of lines crossing a long thin strip fills every box of the strip with points; see \cref{fig:bundle}. The cost of the bundle grows only like the square of a logarithm of the length.

\begin{figure}[ht]
  \centering
  % Generated by make_figures.py -- do not edit by hand.
% seed=4, r=1.0, m=8, K=5, lambda=3.0
\begin{tikzpicture}[x=2.3cm,y=2.3cm]
  \draw[gray!60] (0.0000,0) rectangle (0.5000,0.5000);
  \draw[gray!60] (0.5000,0) rectangle (1.0000,0.5000);
  \draw[gray!60] (1.0000,0) rectangle (1.5000,0.5000);
  \draw[gray!60] (1.5000,0) rectangle (2.0000,0.5000);
  \draw[gray!60] (2.0000,0) rectangle (2.5000,0.5000);
  \draw[gray!60] (2.5000,0) rectangle (3.0000,0.5000);
  \draw[gray!60] (3.0000,0) rectangle (3.5000,0.5000);
  \draw[gray!60] (3.5000,0) rectangle (4.0000,0.5000);
  \draw[gray!60] (4.0000,0) rectangle (4.5000,0.5000);
  \draw[deleted, line width=1.6pt] (0,0.1250) -- (0,0.3750);
  \draw[deleted, line width=1.6pt] (4.5000,0.1250) -- (4.5000,0.3750);
  \draw[gray!80, thin] (-0.15,0.1873) -- (4.6500,0.0815);
  \draw[gray!80, thin] (-0.15,0.2269) -- (4.6500,0.1349);
  \draw[gray!80, thin] (-0.15,0.1424) -- (4.6500,0.1162);
  \draw[gray!80, thin] (-0.15,0.3520) -- (4.6500,0.4321);
  \draw[gray!80, thin] (-0.15,0.3186) -- (4.6500,0.2445);
  \fill[gray!70!black] (0.2564,0.1784) circle (0.02);
  \fill[gray!70!black] (0.3643,0.1760) circle (0.02);
  \fill[gray!70!black] (0.4275,0.1746) circle (0.02);
  \fill[gray!70!black] (0.4649,0.1738) circle (0.02);
  \fill[gray!70!black] (0.5453,0.1720) circle (0.02);
  \fill[gray!70!black] (1.4197,0.1527) circle (0.02);
  \fill[gray!70!black] (2.0081,0.1397) circle (0.02);
  \fill[gray!70!black] (2.5557,0.1277) circle (0.02);
  \fill[gray!70!black] (3.0928,0.1158) circle (0.02);
  \fill[gray!70!black] (3.1645,0.1142) circle (0.02);
  \fill[gray!70!black] (3.2881,0.1115) circle (0.02);
  \fill[gray!70!black] (3.6167,0.1043) circle (0.02);
  \fill[gray!70!black] (4.0554,0.0946) circle (0.02);
  \fill[gray!70!black] (0.7068,0.2105) circle (0.02);
  \fill[gray!70!black] (0.7370,0.2099) circle (0.02);
  \fill[gray!70!black] (1.0473,0.2039) circle (0.02);
  \fill[gray!70!black] (1.4185,0.1968) circle (0.02);
  \fill[gray!70!black] (1.6535,0.1923) circle (0.02);
  \fill[gray!70!black] (1.7187,0.1911) circle (0.02);
  \fill[gray!70!black] (1.9326,0.1870) circle (0.02);
  \fill[gray!70!black] (1.9638,0.1864) circle (0.02);
  \fill[gray!70!black] (2.8726,0.1690) circle (0.02);
  \fill[gray!70!black] (3.5412,0.1561) circle (0.02);
  \fill[gray!70!black] (3.8056,0.1511) circle (0.02);
  \fill[gray!70!black] (3.9246,0.1488) circle (0.02);
  \fill[gray!70!black] (0.2832,0.1401) circle (0.02);
  \fill[gray!70!black] (0.9964,0.1362) circle (0.02);
  \fill[gray!70!black] (1.6244,0.1327) circle (0.02);
  \fill[gray!70!black] (1.8611,0.1315) circle (0.02);
  \fill[gray!70!black] (2.0392,0.1305) circle (0.02);
  \fill[gray!70!black] (2.3438,0.1288) circle (0.02);
  \fill[gray!70!black] (2.5317,0.1278) circle (0.02);
  \fill[gray!70!black] (2.5904,0.1275) circle (0.02);
  \fill[gray!70!black] (2.7117,0.1268) circle (0.02);
  \fill[gray!70!black] (3.2699,0.1238) circle (0.02);
  \fill[gray!70!black] (3.2846,0.1237) circle (0.02);
  \fill[gray!70!black] (3.3004,0.1236) circle (0.02);
  \fill[gray!70!black] (3.6285,0.1218) circle (0.02);
  \fill[gray!70!black] (3.7382,0.1212) circle (0.02);
  \fill[gray!70!black] (3.9932,0.1198) circle (0.02);
  \fill[gray!70!black] (4.2056,0.1186) circle (0.02);
  \fill[gray!70!black] (4.3456,0.1179) circle (0.02);
  \fill[gray!70!black] (0.0725,0.3557) circle (0.02);
  \fill[gray!70!black] (0.2500,0.3587) circle (0.02);
  \fill[gray!70!black] (0.3255,0.3599) circle (0.02);
  \fill[gray!70!black] (0.6592,0.3655) circle (0.02);
  \fill[gray!70!black] (0.7670,0.3673) circle (0.02);
  \fill[gray!70!black] (0.9136,0.3697) circle (0.02);
  \fill[gray!70!black] (1.3716,0.3774) circle (0.02);
  \fill[gray!70!black] (1.5004,0.3795) circle (0.02);
  \fill[gray!70!black] (1.7729,0.3841) circle (0.02);
  \fill[gray!70!black] (2.5551,0.3971) circle (0.02);
  \fill[gray!70!black] (2.5905,0.3977) circle (0.02);
  \fill[gray!70!black] (2.6117,0.3981) circle (0.02);
  \fill[gray!70!black] (2.6984,0.3995) circle (0.02);
  \fill[gray!70!black] (3.1812,0.4076) circle (0.02);
  \fill[gray!70!black] (3.4997,0.4129) circle (0.02);
  \fill[gray!70!black] (3.5901,0.4144) circle (0.02);
  \fill[gray!70!black] (3.7241,0.4167) circle (0.02);
  \fill[gray!70!black] (3.7892,0.4177) circle (0.02);
  \fill[gray!70!black] (3.9940,0.4212) circle (0.02);
  \fill[gray!70!black] (4.0086,0.4214) circle (0.02);
  \fill[gray!70!black] (4.4068,0.4280) circle (0.02);
  \fill[gray!70!black] (1.0316,0.3004) circle (0.02);
  \fill[gray!70!black] (1.4741,0.2935) circle (0.02);
  \fill[gray!70!black] (2.5458,0.2770) circle (0.02);
  \fill[gray!70!black] (2.5520,0.2769) circle (0.02);
  \fill[gray!70!black] (2.6656,0.2751) circle (0.02);
  \fill[gray!70!black] (3.7927,0.2577) circle (0.02);
  \fill[gray!70!black] (4.2902,0.2500) circle (0.02);
  \fill[gray!70!black] (4.4663,0.2473) circle (0.02);
  \draw[retained, thick] (0.2564,0.1784) -- (0.7068,0.2105) -- (1.4741,0.2935) -- (1.6535,0.1923) -- (2.0081,0.1397) -- (2.6656,0.2751) -- (3.2699,0.1238) -- (3.7927,0.2577) -- (4.2902,0.2500);
  \fill[retained] (0.2564,0.1784) circle (0.03);
  \fill[retained] (0.7068,0.2105) circle (0.03);
  \fill[retained] (1.4741,0.2935) circle (0.03);
  \fill[retained] (1.6535,0.1923) circle (0.03);
  \fill[retained] (2.0081,0.1397) circle (0.03);
  \fill[retained] (2.6656,0.2751) circle (0.03);
  \fill[retained] (3.2699,0.1238) circle (0.03);
  \fill[retained] (3.7927,0.2577) circle (0.03);
  \fill[retained] (4.2902,0.2500) circle (0.03);
  \draw[|-|] (0,-0.12) -- (0.5000,-0.12) node[midway, below, font=\small] {$r/2$};
  \node[font=\small, anchor=east] at (-0.17,0.2500) {$R$};
\end{tikzpicture}

  \caption{The mechanism behind \cref{prop:noexp}, numerical illustration with $r=1$, $m=8$, $K=5$ lines and linear point intensity~$3$. The lines pass through the middle halves (thick amber) of the two vertical sides of the strip $R=[0,(m+1)r/2]\times[0,r/2]$ with slopes at most $a/(4d)$, as in the proof of \cref{lem:strip}. They meet each of the $m+1$ boxes of side $r/2$ in a segment of length at least $r/2$, so each box receives Poisson points of mean at least $\lambda Kr/2$. One point per box is highlighted; consecutive highlighted points are at distance less than~$2r$, so their disks form a connected chain across~$R$. Only the strip is shown, and the vertical scale equals the horizontal one.}
  \label{fig:bundle}
\end{figure}

\begin{proposition}[Quasi-polynomial lower bound]
\label{prop:noexp}
Let $\lambda,\mu,r>0$, and for $m\in\N$ let $C_m$ be the event that a connected component of~$\Xi_\lambda$ meets both $[0,r/2]^2$ and $[mr/2,(m+1)r/2]\times[0,r/2]$. Then
\begin{equation*}
  \liminf_{m\to\infty}\frac{\log\P(C_m)}{(\log m)^2}\ge-\frac{4}{\lambda r}\,.
\end{equation*}
In particular, $\P(C_m)\ge\exp(-m^\delta)$ for every $\delta>0$ and all large~$m$.
\end{proposition}

\begin{proof}
Let $R\coloneqq[0,(m+1)r/2]\times[0,r/2]$ and $Q'_k\coloneqq[kr/2,(k+1)r/2)\times[0,r/2]$ for $0\le k\le m$. Let $M$ be the number of lines of~$\cL$ meeting both vertical sides of~$R$. By \cref{lem:strip} with $a=r/2$ and $d=(m+1)r/2$, the variable~$M$ is Poisson with mean $\nu_m\in[\mu r/(10\pi(m+1)),\mu r/\pi]$, and every such line meets every~$Q'_k$ in a segment of length at least~$r/2$. Hence $\Lambda(Q'_k)\ge Mr/2$ for all~$k$.

Conditionally on~$\cL$, the numbers of points in the disjoint boxes~$Q'_k$ are independent, and each box contains a point with probability at least $1-\e^{-\lambda rM/2}$. If every box contains a point $x_k\in Q'_k$, then $\lVert x_k-x_{k+1}\rVert_2\le(\sqrt5/2)r<2r$, so the disks $B(x_k,r)$ form a connected set meeting both end boxes, and $C_m$ occurs. Put $j\coloneqq\lceil4\log(m+1)/(\lambda r)\rceil$. On $\{M\ge j\}$ we have $\e^{-\lambda rM/2}\le(m+1)^{-2}$, and Bernoulli's inequality gives
\begin{equation*}
  \P(C_m)\ge\P(M\ge j)\bigl(1-(m+1)^{-2}\bigr)^{m+1}\ge\frac12\,\e^{-\nu_m}\frac{\nu_m^j}{j!}\ge\frac12\,\e^{-\mu r/\pi}\Bigl(\frac{\mu r}{10\pi(m+1)j}\Bigr)^j\,,
\end{equation*}
where we used $j!\le j^j$. Since $j=(4+o(1))\log m/(\lambda r)$, the logarithm of the right-hand side is $-(4+o(1))(\log m)^2/(\lambda r)$.
\end{proof}

The bound holds for every $\lambda>0$. Hence no argument that yields exponential decay of connection probabilities can prove a subcritical phase; any proof has to accommodate decay of quasi-polynomial order. We expect that the order $(\log m)^2$ is sharp in the subcritical regime, but we do not prove a matching upper bound.

% =================================================================
\section{A quenched criterion}
\label{sec:criterion}

By \cref{prop:quenched}, it suffices to find, for a positive proportion of line configurations, one intensity without percolation. Conditionally on the lines, occupied boxes are independent, so a union bound over lattice paths only involves the line counts of the boxes. The cost of this union bound is governed by a growth rate of the line configuration.

For $z\in\Z^2$ let $Q(z)\coloneqq rz+[0,r]^2$ and let $N(z)$ be the number of lines of~$\cL$ meeting~$Q(z)$. A \emph{$2$-path} of length~$n$ from~$z_0$ is a sequence $\gamma=(z_0,\dots,z_n)$ of distinct points of~$\Z^2$ with $\lVert z_i-z_{i-1}\rVert_\infty\le2$; let $\Gamma_n(z_0)$ denote the set of these paths, so that $\#\Gamma_n(z_0)\le24^n$. Define
\begin{equation*}
  Y_n(z_0)\coloneqq\sum_{\gamma\in\Gamma_n(z_0)}\ \prod_{i=0}^nN(z_i),\qquad Y_n^*\coloneqq\max_{\lVert z_0\rVert_\infty\le n}Y_n(z_0),\qquad g^*\coloneqq\limsup_{n\to\infty}\frac1n\log Y_n^*\,,
\end{equation*}
with $\log0=-\infty$. The quantity $g^*\in[-\infty,\infty]$ is a function of~$\cL$ alone.

\begin{theorem}[Quenched path criterion]
\label{thm:criterion}
For every realization of~$\cL$ meeting every bounded set in finitely many lines, $\lambda_c(\cL)\ge\e^{-g^*}/(\sqrt2\,r)$. Consequently, if $\P(g^*<\infty)>0$, then $\lambda_c(\mu,r)>0$.
\end{theorem}

\begin{proof}
If $g^*=\infty$, the lower bound is zero and needs no argument. Otherwise fix $\lambda>0$ with $\lambda\sqrt2\,r\,\e^{g^*}<1$, and choose a finite $a>g^*$ with $\lambda\sqrt2\,r\,\e^a<1$. Such an $a$ exists also when $g^*=-\infty$. Call $z$ occupied if $X_\lambda(Q^\circ(z))\ge1$, where $Q^\circ(z)\coloneqq rz+[0,r)^2$. The boxes $Q^\circ(z)$ are disjoint, so conditionally on~$\cL$ the events $\{z\text{ occupied}\}$ are independent. Each line meets $Q^\circ(z)$ in a segment of length at most $\sqrt2\,r$, so
\begin{equation*}
  \P(z\text{ occupied}\,\vert\,\cL)=1-\e^{-\lambda\Lambda(Q^\circ(z))}\le\lambda\sqrt2\,r\,N(z)\,.
\end{equation*}

Let $A_{z_0}$ be the event that an unbounded component of~$\Xi_\lambda$ contains the disk of a Cox point in~$Q^\circ(z_0)$. On this event, the Cox points whose disks belong to that component are infinitely many, and two of its points $x\in Q^\circ(z)$ and $y\in Q^\circ(w)$ with intersecting disks satisfy $\lvert x-y\rvert<2r$, hence $\lVert z-w\rVert_\infty\le2$. The occupied sites containing its points therefore form an infinite set containing~$z_0$ that is connected for $2$-adjacency. Since this graph is locally finite, it contains a vertex at graph distance~$n$ from~$z_0$ for every~$n$, and a shortest path to it is a $2$-path in $\Gamma_n(z_0)$ whose sites are all occupied. By the union bound and conditional independence, for $n\ge\lVert z_0\rVert_\infty$,
\begin{equation*}
  \P(A_{z_0}\,\vert\,\cL)\le\sum_{\gamma\in\Gamma_n(z_0)}\prod_{i=0}^n\lambda\sqrt2\,r\,N(z_i)\le(\lambda\sqrt2\,r)^{n+1}\,Y_n^*\,.
\end{equation*}
For all large~$n$ we have $Y_n^*\le\e^{an}$, so the right-hand side tends to zero. A union over $z_0\in\Z^2$ gives $\theta_\cL(\lambda)=0$, which proves the first assertion. If $\P(g^*<\infty)>0$, then $\lambda_c(\cL)>0$ with positive probability, and \cref{prop:quenched} gives $\lambda_c(\mu,r)>0$.
\end{proof}

The criterion involves no point process and no restriction on~$\mu r$. Since $\#\Gamma_n(z_0)\le24^n$, it reduces to a greedy-path statement for the field $\log N$:
\begin{equation*}
  Y_n(z_0)\le24^n\max_{\gamma\in\Gamma_n(z_0)}\ \prod_{i=0}^nN(z_i)\,.
\end{equation*}

\begin{corollary}[Greedy-path criterion]
\label{cor:greedy}
Let
\begin{equation*}
  h^*\coloneqq\limsup_{n\to\infty}\frac1n\max_{\lVert z_0\rVert_\infty\le n}\ \max_{\gamma\in\Gamma_n(z_0)}\ \sum_{i=0}^n\log N(z_i)\,.
\end{equation*}
Then $g^*\le\log24+h^*$. In particular, if $\P(h^*<\infty)>0$, then $\lambda_c(\mu,r)\ge\e^{-\log24-\operatorname{ess\,inf}h^*}/(\sqrt2\,r)>0$.
\end{corollary}

% =================================================================
\section{Direct annealed bounds and dependence}
\label{sec:obstacles}

The natural way to prove $g^*<\infty$ is a first-moment bound on~$Y_n$ followed by Borel--Cantelli. The same bundles as in \cref{prop:noexp} make this impossible.

\begin{proposition}[Superexponential annealed moments]
\label{prop:annealed}
For every $\mu,r>0$,
\begin{equation*}
  \lim_{n\to\infty}\frac1n\log\E Y_n(0)=\infty\,.
\end{equation*}
\end{proposition}

\begin{proof}
The straight path $(0,e_1,\dots,ne_1)$ belongs to~$\Gamma_n(0)$. Let $M$ be the number of lines meeting both vertical sides of $[0,(n+1)r]\times[0,r]$. By \cref{lem:strip}, $M$ is Poisson with mean $\nu_n\in[\mu r/(5\pi(n+1)),2\mu r/\pi]$, and every such line meets every $Q(ke_1)$, so $N(ke_1)\ge M$ for $0\le k\le n$. With $j\coloneqq\lceil n/\log n\rceil$,
\begin{equation*}
  \E Y_n(0)\ge\E M^{n+1}\ge j^{n+1}\,\e^{-\nu_n}\frac{\nu_n^j}{j!}\ge\e^{-2\mu r/\pi}\,j^{n+1}\Bigl(\frac{\mu r}{5\pi(n+1)j}\Bigr)^j\,,
\end{equation*}
and the logarithm of the right-hand side is $n\log n-O(n\log\log n)$.
\end{proof}

The bound comes from events of probability about $\e^{-2n}$, in which roughly $n/\log n$ lines run along a fixed straight path, so it does not contradict $g^*<\infty$. It shows only that the union bound over paths must be organised around the line configuration rather than taken in expectation. The following remarks explain why other standard routes fail as well. They are heuristic and are not used elsewhere.

\paragraph{Decoupling.} In a multiscale scheme with scales $L_{n+1}=aL_n$, two windows of size~$L_n$ at distance $aL_n$ are both met by a Poisson number of lines with mean of order $\mu L_n/a$, which grows with the scale. Conditionally on all lines, the point configurations in the two windows are independent, but occupied crossing events increase with the lines, so the shared lines correlate them positively. The additive decoupling of \cref{lem:mixing} has error of order $\mu L_n/a$ and is useless. The deletion of directions in~\cite{HJ26} removes all shared lines but decreases occupied events, which is the wrong direction for an upper bound~\cite[Remark~6.1]{HJ26}.

\paragraph{Sprinkling.} For increasing events one would like to dominate the shared lines in each window by additional independent lines. The shared lines have angular density $\mu/\pi$ on an interval of length of order~$1/a$; an independent sprinkle of intensity $\varepsilon\mu$ has density $\varepsilon\mu/\pi$ there and cannot dominate them unless $\varepsilon\ge1$.

\paragraph{Vacant circuits.} The complement of~$\Xi_\lambda$ contains every point at distance more than~$r$ from all lines, that is, the inner parallel bodies of the cells of the line tessellation. A strip around a line can be crossed at a point $p$ of an edge if no other line comes within distance~$3r$ of~$p$ and no Cox point lies within distance~$2r$ of~$p$. A planar duality argument then shows that $\Xi_\lambda$ does not percolate if the edges without such a crossing point have only finite clusters in the line arrangement. However, two lines at angle~$\alpha$ stay within distance~$3r$ of each other over a length of order $r/\alpha$, and since the angle at a typical intersection has density proportional to $\sin\alpha$, these blocked stretches have length~$L$ with $\E L^2=\infty$ (logarithmically). A first-moment count of blocked paths in the arrangement then diverges for every $\mu>0$, in analogy with Boolean models with random radii, where positivity of the critical intensity is tied to a moment condition on the radius~\cite{Gou08}. Crossing two strips at once removes this particular obstacle but leaves bundles of three or more lines, and the duality becomes considerably more involved.

% =================================================================
\section{Towards the greedy-path estimate}
\label{sec:greedy}

\Cref{cor:greedy} asks for linear growth of the maximal sum of $\log N$ along paths, uniformly over starting points in a box of size~$n$. For independent weights this is the theory of greedy lattice paths and animals: in dimension two, the maximum over paths of length~$n$ grows linearly under the condition $\int_0^\infty\P(W>w)^{1/2}\dd w<\infty$ on the weight~$W$~\cite{CGGK93,Mar02}. The field $\log N$ has superexponential one-point tails, but it is dependent at all distances through shared lines. The following lemma controls the shared lines quenched.

\begin{lemma}[Shared lines between distant boxes]
\label{lem:shared}
For every $k\in\N$ there is $c_k=c_k(\mu r)<\infty$ such that, almost surely, for all large~$n$, any two boxes $Q(z),Q(w)$ with $\lVert z\rVert_\infty,\lVert w\rVert_\infty\le n$ and $\lVert z-w\rVert_2\ge c_kn^{6/k}$ are both met by fewer than~$k$ lines. Moreover, almost surely $\max_{\lVert z\rVert_\infty\le n}N(z)\le\log n$ for all large~$n$.
\end{lemma}

\begin{proof}
Each box lies in a disk of radius $r/\sqrt2$ around its center, so by \cref{lem:mixing} the lines meeting two boxes at distance $D=r\lVert z-w\rVert_2\ge4r$ have measure at most $2\mu r^2/D$. Their number is Poisson, and $\P(\mathrm{Poisson}(\nu)\ge k)\le\nu^k/k!\le(\e\nu/k)^k$. There are at most $(2n+1)^4$ pairs, so the probability of a violation at stage~$n$ is at most $(2n+1)^4(2\e\mu r/(k\lVert z-w\rVert_2))^k\le81\,n^{-2}$ if $c_k\coloneqq\max\{4,2\e\mu r/k\}$, using $\lVert z-w\rVert_2^{-k}\le c_k^{-k}n^{-6}$. The Borel--Cantelli lemma gives the first claim. For the second, $N(z)$ is Poisson with mean $4\mu r/\pi$, and $(2n+1)^2(4\e\mu r/(\pi\log n))^{\log n}$ is summable.
\end{proof}

\begin{conjecture}[Linear greedy growth]
\label{conj:greedy}
For every $\mu,r>0$, almost surely $h^*<\infty$.
\end{conjecture}

By \cref{cor:greedy}, \cref{conj:greedy} implies $\lambda_c(\mu,r)>0$ for all $\mu,r>0$.

\paragraph{Heuristic support.} Large values of $N$ along a path arise in two ways. Many lines can cross near a single box at different angles; these clusters are local and behave like independent weights with superexponential tails, which the greedy theory handles. Alternatively, $j$ nearly parallel lines can form a bundle along a stretch of length~$t$. By the computation in the proof of \cref{lem:shared}, such bundles occur near a given point with probability of order $(\mu r^2/t)^j/j!$, and a path following one collects about $t\log j$. Treating bundles as rewards attached to sites, the tail of the reward is of order $t^{-j}$ for each~$j\ge2$, and $\int\P(W>w)^{1/2}\dd w$ converges once only bundles with $j\ge3$ contribute substantially; bundles with $j\le2$ contribute at most $\log2$ per site. This suggests linear growth. The rigorous difficulty is that bundles are built from shared lines and are therefore neither independent nor of bounded range, and that the union bound over paths has to exploit the low entropy of paths that follow a bundle, since \cref{prop:annealed} excludes an expectation bound.

\paragraph{Possible route.} For $N\ge1$, one may write $\log N\le(1+\sum_{k\ge1}\1\{N\ge2^k\})\log 2$ and prove, for each level set $D_m\coloneqq\{z\colon N(z)\ge m\}$, that every path in $\Gamma_n(z_0)$ meets $D_m$ in at most $\varepsilon_mn$ sites, with $\sum_k\varepsilon_{2^k}<\infty$. By \cref{lem:shared} only $m\le\log n$ matter. For a fixed path, the event that a given set of its sites lies in~$D_m$ can be expanded according to which lines explain the large counts. Lines whose footprint on the path is short form a Poisson process with bounded incidences and give exponentially small probabilities; long footprints force the path to follow a line over a long stretch, which reduces the number of admissible paths. Making this trade-off between probability and path entropy precise is the missing step.

\paragraph{Expected size of the threshold.} We expect, heuristically, that
\begin{equation*}
  \lambda_c(\mu,r)\ge(1-o(1))\,\frac{\log(1/(\mu r))}{2r}\qquad\text{as }\mu r\to0\,.
\end{equation*} To percolate, occupied stretches must run along lines between intersections, which occur at rate $2\mu/\pi$, while gaps longer than~$2r$ between points occur at rate $\lambda\e^{-2\lambda r}$; at most three continuations at each intersection suggest that percolation fails once $\lambda\e^{-2\lambda r}>4\mu/\pi$. Together with~\cite[Theorem~2.1]{HJ26}, this would determine the order $r^{-1}\log(1/(\mu r))$ of~$\lambda_c$ for sparse lines. For dense lines, $\mu r\to\infty$, the Cox process resembles a Poisson process of intensity~$\lambda\mu$, and we expect $\lambda_c(\mu,r)$ of order $1/(\mu r^2)$; in this regime the upper bound of order $1/r$ in~\cite{HJ26} is not sharp.

% =================================================================
\begin{samepage}
\section*{Acknowledgements}

BJ acknowledges the financial support of the Leibniz Association within the Leibniz Junior Research Group on Probabilistic Methods for Dynamic Communication Networks (DYCOMNET) as part of the Leibniz Competition (grant no.~J105/2020), of the Berlin Mathematics Research Center MATH+ through the project EF-MA-Sys-2 on Information Flow \& Emergent Behavior in Complex Networks, and of the Deutsche Forschungsgemeinschaft (DFG) through Project P27 within the SPP~2265.

\section*{Use of AI tools}

The simulations, proofs, computations, figures and the text of this note were developed with substantial assistance of the AI systems Claude (Anthropic) and ChatGPT Astra~6 (OpenAI), in an interactive process directed by the authors.

\end{samepage}

\bibliographystyle{amsalpha}
\bibliography{references}

\end{document}
